% ============================================================
\chapter{KẾT LUẬN VÀ HƯỚNG PHÁT TRIỂN}
\label{ch:ket_luan}
% ============================================================

\section{Kết luận}

Báo cáo đã nghiên cứu một cách hệ thống hai chiến lược xử lý mất cân bằng lớp (cấu hình Exp 2 với Cross-Entropy trọng số ENet cố định và cấu hình Exp 4 với tổ hợp Focal+Dice+CE học trọng số thích nghi) trên bốn kiến trúc phân vùng ngữ nghĩa ảnh lâm nghiệp UAV phổ biến là DeepLabV3+, SegFormer, UNet và UNet++. 

Các kết luận khoa học chính rút ra từ nghiên cứu bao gồm:
\begin{itemize}[leftmargin=1.4cm]
  \item \textbf{Loss thích nghi (Exp 4) vượt trội toàn diện:} Việc sử dụng tổ hợp hàm mất mát thích nghi thông qua học trọng số đa nhiệm (Uncertainty Weighting) giúp tăng vượt trội điểm mIoU kiểm thử từ $+4{,}15$ pp đến $+4{,}88$ pp trên tất cả bốn kiến trúc mạng, đạt hiệu năng cao nhất trên kiến trúc UNet++ với mIoU bằng $0{,}8374$.
  \item \textbf{Cải thiện mạnh mẽ chất lượng phân vùng lớp hiếm:} Exp 4 mang lại sự gia tăng đột phá về IoU trên các lớp đối tượng đặc biệt hiếm và khó nhận diện như cây đổ (\textit{Fallen trees}), hàng rào (\textit{Fence}) và ô tô (\textit{Car}), đồng thời bảo toàn hoặc tối ưu hóa hiệu năng phân vùng trên các lớp đa số (không bị suy giảm như trong Exp 2).
  \item \textbf{Sự đánh đổi về độ bao phủ trên vật thể nhỏ:} Thực nghiệm bộc lộ một đặc tính quan trọng của loss thích nghi trong Exp 4 là cải thiện rất mạnh độ chính xác hình học (IoU) nhưng lại làm giảm nhẹ độ bao phủ Recall (tăng tỷ lệ bỏ sót Miss Rate đối với các vật thể nhỏ đơn lẻ từ $3\%$ đến $8\%$). Đây là một thông tin thiết kế quan trọng cho các kỹ sư khi triển khai mô hình trong thực tế.
\end{itemize}

\section{Hướng phát triển}
\label{sec:huong_morong}

Trong tương lai, nghiên cứu có thể được mở rộng theo các hướng sau:
\begin{enumerate}[leftmargin=1.4cm]
  \item \textbf{Khắc phục điểm yếu bỏ sót vật thể nhỏ:} Nghiên cứu kết hợp các kỹ thuật tăng cường dữ liệu chuyên biệt cho vật thể nhỏ (như Copy-Paste Augmentation) hoặc tích hợp thêm các hàm mất mát hướng ranh giới (như Boundary Loss) vào tổ hợp thích nghi để giảm tỷ lệ bỏ sót (Miss Rate) trong Exp 4.
  \item \textbf{Đánh giá chuyển miền thực tế:} Thực hiện kiểm chứng tính thực tiễn của mô hình đã huấn luyện bằng cách ánh xạ nhãn và đánh giá khả năng phân vùng trên tập dữ liệu ảnh chụp thực tế ngoài đời thực (như WildUAV), đồng thời nghiên cứu các kỹ thuật thích nghi miền không giám sát (Unsupervised Domain Adaptation).
  \item \textbf{Hậu xử lý và tối ưu hóa biên:} Ứng dụng các phương pháp hậu xử lý như CRF (Conditional Random Fields) hoặc kỹ thuật suy diễn đa tỷ lệ khi kiểm thử (Test-Time Augmentation - TTA) để làm mịn ranh giới phân vùng các đối tượng dạng đường biên mảnh như hàng rào.
  \item \textbf{Tối ưu hóa triển khai trên thiết bị biên:} Thử nghiệm chuyển đổi mô hình sang các định dạng tối ưu (như TensorRT, ONNX) và đánh giá tốc độ xử lý (FPS), điện năng tiêu thụ trên các thiết bị máy tính nhúng của UAV (như NVIDIA Jetson) để hướng tới phân đoạn thời gian thực.
\end{enumerate}
